\section{Conclusion}

Return to the difficulty of manually constructing reusable BDI plan
libraries as the knowledge-engineering problem addressed by the paper.

Conclude that reusable behaviour obtained across planning instances can
be transformed into a parameterised BDI plan library that supports
achievement goals beyond the instances used to construct it.

Conclude that the same domain-level library can be retained and reused
for subsequent user queries rather than generating domain behaviour
again for each new goal.

Conclude that the query interface supports both achievement and
temporally extended goals expressed through controlled natural language.

Relate the experimental results to the two main claims by showing that
completing the reusable library improves held-out goal achievement and
that query-specific coordination enables the generated library to
realise the supported temporal goals.

State that the current guarantees apply to the supported planning
fragment, generated plan space, and accepted query compositions rather
than to arbitrary planning domains or temporal specifications.

Identify the availability of representative training instances as a
current assumption of domain-level library generation.

Introduce parameterised PDDL problem generators as a way to
automatically produce reproducible training instances for new domains
while reducing the remaining engineering effort required for
domain-level library generation.
